% BEGIN PROFESSOR FEEDBACK PF11: Method - Temporal motif extraction
\section{Temporal motif extraction}
\label{sec:motif_extraction}

The baseline describes an address by totals, averages, and timing summaries.
It does not retain the full order of transfers between particular partners.
Two addresses can have the same number of
transfers, the same partners, and the same active span, yet behave very
differently: one might receive money from many senders within minutes and pass
it straight on, while the other collects the same money slowly over months.
Candidate temporal patterns capture these differences. Recall from
Section~\ref{sec:motifs} that a candidate is a short, time-ordered sequence of
transfers that all fall within a window of length $\delta$. Choosing a focal
address determines its local role in an occurrence. Stacking the address's
local role counts into a vector gives its motif fingerprint
(Equation~\ref{eq:motif_fingerprint}). This section describes how those
fingerprints are counted and turned into features.

\subsection{Subgraph counting and retained local roles}
\label{sec:motifcountingandretainedroles}

The counting itself is done by Raphtory, whose
\texttt{local\_temporal\_three\_node\_motifs} routine is written in Rust and
processes addresses in parallel. For every address it returns a \emph{local}
count for each motif role, that is, the number of candidate occurrences in
which the address plays that particular part. A star pattern is counted only for
the address at its centre, and a two-node pattern for both of the addresses
involved.

Raphtory reports 40 such roles per address: 24 star roles, 8 two-node roles, and
8 triangle roles. The star roles comprise three classes of eight in-and-out
direction patterns each. The four global two-node patterns each have two
endpoint roles, giving eight local two-node counts. This accounts for the
difference between 36 global patterns and 40 local role features. The output
is not a catalogue of 36 global motifs with four additional patterns.
The global patterns are shown in
Figure~\ref{fig:motifs_paranjape}. The final HuggingFace experiment retains all
40 roles because its released-neighbourhood closure supports triangle counting.
It also evaluates the 32 star and two-node roles separately to measure whether
triangles add signal. Real-CATS retains those 32 roles per layer because its
collected histories cannot guarantee the edge between two partners. An address
that takes part in no retained candidate pattern keeps a fingerprint of all zeros rather
than being dropped from the study.

The retained roles fall into a few structural classes, each with a
straightforward reading in a transaction network.
Table~\ref{tab:motif_classes} lists them together with the kind of coordinated
behaviour they could represent. These interpretations do not label an occurrence
as fraudulent. In every example, the timing condition is
$t_3-t_1 \leq \delta$, with $\delta$ drawn from the five-minute to 30-day grid
in Section~\ref{sec:temporal_window_sweep}. There is no separate threshold for
\emph{rapid} activity, and an occurrence admitted at 30 days need not be rapid in
an everyday sense.

\begin{table}[h]
\centering
\caption{Retained structural motif classes (after \citet{paranjape2017motifs})
and their behavioural interpretation in Ethereum transaction networks. Each
class covers several $\delta$-temporal roles that differ in the direction and
order of the three transfers.}
\label{tab:motif_classes}
\begin{tabular}{p{3.5cm} p{4.5cm} p{5cm}}
\toprule
\textbf{Class} & \textbf{Ordered transfers} & \textbf{Possible interpretation} \\
\midrule
Two-node, same direction &
    $A \to B,\; A \to B,\; A \to B$ &
    Repeated transfers, which may indicate instalment-based extraction \\[4pt]
Two-node, mixed direction &
    $A \to B,\; B \to A,\; A \to B$ &
    Back-and-forth transfers within $\delta$, as in wash-trading or fee manipulation \\[4pt]
Star, all incoming &
    $B \to A,\; C \to A,\; C \to A$ &
    Funds consolidating at a single hub, as in money-mule collection \\[4pt]
Star, all outgoing &
    $A \to B,\; A \to C,\; A \to C$ &
    Mass distribution from one source, as in airdrop farming or Ponzi payouts \\[4pt]
Star, mixed &
    $A \to B,\; C \to A,\; A \to C$ &
    A relay node forwarding funds, as in a layering intermediary \\[4pt]
Triangle &
    $A \to B,\; B \to C,\; C \to A$ &
    A closed transfer sequence, shown here as a cycle. Other triangle roles
    have different directions. \\
\bottomrule
\end{tabular}
\end{table}

The terminology in Table~\ref{tab:motif_terminology} applies throughout this
chapter. Each feature column describes a local role, each class groups related
roles, and the fingerprint contains the complete set of retained counts for
one address.

\subsection{Ordering transfers with equal timestamps}

Because motifs are ordered sequences, the transfers have to be put in a definite
order before they can be counted. Time supplies most of that order, but many
transfers carry the same timestamp because they sit in the same block, and time
alone cannot separate them. When the data records where a transfer sits within its
block, that finer order is used. The HuggingFace data does not record this, so its
transfers that share a timestamp have no fixed order. To make sure the results do
not rest on an arbitrary choice, the whole analysis is run under seven versions of
the order. All orderings preserve the order of transfers with distinct
timestamps. They differ only in the order assigned to transfers sharing the
same timestamp: one fixed order, its reversal within each timestamp group, and
five random shuffles within those groups. These differences can change which
candidate temporal pattern a sequence matches and which local roles are counted.
Section~\ref{sec:evaluation_protocol} reports whether the conclusions hold under
all seven. The order is computed without loading the whole table into memory, and
the sorted transfers are passed to Raphtory in large batches.

\subsection{Fingerprint representations}

The retained counts are stored under stable names such as \texttt{star\_pre\_iii},
\texttt{two\_node\_ioi}, and \texttt{triangle\_1}. For RQ1 and RQ2 each count
$c$ is replaced by $\log(1+c)$. Some highly active addresses have much larger
counts than most targets. On the raw scale these values would dominate the model.
Taking the logarithm compresses this range, and adding one first means that a
count of zero stays zero. The final HuggingFace detection step uses all 40
transformed values, either on their own or joined to the baseline of
Section~\ref{sec:graph_construction}. A parallel 32-role representation and a
triangles-only representation support the ablation analysis. How the fingerprint
and baseline are combined into the RQ1 feature sets is described in
Section~\ref{sec:detection_methods}.

The representation comparisons also use \emph{motif volume}, defined here as
the sum of an address's retained local role counts. This is a local summary,
not the number of distinct global pattern types.

RQ3 and RQ4 keep the five Real-CATS layers separate. Because the layers differ greatly
in how many transfers they contain, comparing raw counts would let the busiest
layer dominate. Instead, within each layer the counts are turned into proportions by dividing each
count by the layer total, and the square root of each proportion is taken. This is
known as a Hellinger transform \citep{legendre2001transformations}:
\[
    h_{v,j,l}^{\delta}
      = \sqrt{\frac{c_{v,j,l}^{\delta}}{\sum_{k=1}^{32}c_{v,k,l}^{\delta}}},
\]
where $c_{v,j,l}^{\delta}$ is the count of local role $j$ for address $v$ in
layer $l$ at window $\delta$. The expression applies when the layer's total
role count is positive. Using proportions removes the total count, so only
the mix of local roles remains, and
the square root keeps a few large proportions from dominating, much as the
logarithm does for the counts. A layer with no retained candidate occurrences is
kept as all zeros, including a layer that has transfers but no complete
three-edge pattern. The
same counts in log form, as used for RQ1 and RQ2, are also kept as a check.
Concatenating the five layers gives 160 local role features for each Real-CATS
address, with 32 components per layer.

For every window the pipeline additionally records whether each address has any
nonzero count at all. This coverage measure shows how the availability of
fingerprints changes with the window. Reporting coverage alongside performance
helps assess their relationship, but does not isolate whether a score increase
comes from lower sparsity or a more informative time scale.

\subsection{Temporal-window sweep}
\label{sec:temporal_window_sweep}

The counting is repeated over a fixed grid of nine windows,
\[
    \delta \in \{300,\,1800,\,3600,\,21600,\,86400,\,259200,\,604800,\,
                  1209600,\,2592000\}
    \text{ seconds},
\]
which correspond to five minutes, 30 minutes, one hour, six hours, one day, three
days, one week, two weeks, and 30 days. All windows use the same addresses and
the same train and test split, and the window is chosen only on the development
data, as described in Section~\ref{sec:evaluation_protocol}. RQ2 studies how
detection changes across this grid, and RQ4 asks whether separate windows for
each category improve classification over one shared window. A selected window
is the maximum span allowed for a three-edge candidate motif. It is not an estimate of the
total duration of a fraud campaign.
% END PROFESSOR FEEDBACK PF11
